\documentclass[11pt,a4paper]{article}
\usepackage{amsmath,amssymb,amsthm,mathtools}
\usepackage[margin=2.5cm]{geometry}
\usepackage{booktabs}
\usepackage{hyperref}

\theoremstyle{plain}
\newtheorem{theorem}{Theorem}
\newtheorem{conjecture}[theorem]{Conjecture}
\newtheorem{proposition}[theorem]{Proposition}
\newtheorem{corollary}[theorem]{Corollary}
\newtheorem{lemma}[theorem]{Lemma}
\theoremstyle{remark}
\newtheorem{remark}[theorem]{Remark}

\DeclareMathOperator{\CF}{CF}

\title{Two Parametric Families of Polynomial Continued Fractions\\
for Reciprocals of Logarithms and Multiples of $1/\pi$}
\author{}
\date{April 2026}

\begin{document}
\maketitle

\begin{abstract}
We present two infinite parametric families of polynomial continued fractions
(PCFs) discovered through systematic modular signature sweeps:
a \emph{logarithmic ladder} producing $1/\ln(k/(k{-}1))$ for all real $k > 1$,
and a \emph{pi family} producing $2^{2m+1}/(\pi \binom{2m}{m})$ for all
non-negative integers~$m$.
Both families are certified to over $1\,200$ decimal digits using Arb ball
arithmetic with rigorous consecutive-convergent bracketing.
For the logarithmic ladder, we give a \textbf{complete proof} by exhibiting
closed forms $p_n = (n{+}1)!\,k^{n+1}$ and $q_n = (n{+}1)!\sum_{j=0}^n k^{n-j}/(j{+}1)$
for the convergent numerators and denominators, reducing the limit to
the Taylor series $\sum x^j/(j{+}1) = -\ln(1{-}x)/x$.
For the pi family, the base case $m=0$ is \textbf{proved} via
$p_n = (2n{+}1)!!$ and $q_n = (2n{+}1)!!\sum_{j=0}^n j!/(2j{+}1)!!$,
with the limit given by the classical identity
$\pi/2 = \sum j!/(2j{+}1)!!$; the full family then follows by an
explicit binomial recurrence.
\end{abstract}

% ──────────────────────────────────────────────────────────────────────
\section{Introduction}
% ──────────────────────────────────────────────────────────────────────

A polynomial continued fraction (PCF) is an expression of the form
\begin{equation}\label{eq:pcf}
  b_0 + \cfrac{a_1}{b_1 + \cfrac{a_2}{b_2 + \cfrac{a_3}{b_3 + \cdots}}}
\end{equation}
where $a_n = \alpha(n)$ and $b_n = \beta(n)$ are polynomials in~$n$.
The study of such continued fractions has a rich history going back to
Euler, Gauss, and Ramanujan.
Recent algorithmic approaches---notably the Ramanujan Machine
project~\cite{raayoni2021,elimelech2024}---have revived interest in
systematic discovery of new PCF identities for fundamental constants.

In this note we report two infinite parametric families of PCFs
discovered through a \emph{modular signature sweep}: fixing the
denominator polynomial $b(n) = dn + e$ for small integers $d$ and~$e$,
then exhaustively searching quadratic numerator polynomials~$a(n)$.
The key finding is that the $d=3$ family produces both $\pi$- and
$\ln$-related constants, complementing the classical $d=2$ family
(Brouncker's $4/\pi$, Euler's continued fraction for~$e$).

% ──────────────────────────────────────────────────────────────────────
\section{Main Results}
% ──────────────────────────────────────────────────────────────────────

\subsection{The Logarithmic Ladder}

\begin{theorem}[Logarithmic Ladder]\label{thm:log}
For all real $k > 1$, the polynomial continued fraction with
\[
  a(n) = -k\,n^2, \qquad b(n) = (k+1)\,n + k
\]
converges to $\dfrac{1}{\ln(k/(k-1))}$.
Explicitly:
\[
  k + \cfrac{-k}{(2k+1) + \cfrac{-4k}{(3k+2) + \cfrac{-9k}{(4k+3) + \cdots}}}
  \;=\; \frac{1}{\ln\!\bigl(\frac{k}{k-1}\bigr)}.
\]
\end{theorem}

\begin{proof}
We show by induction that the convergent numerators and denominators satisfy
\begin{equation}\label{eq:pn}
  p_n = (n+1)!\, k^{n+1}, \qquad
  q_n = (n+1)!\, \sum_{j=0}^{n} \frac{k^{n-j}}{j+1}.
\end{equation}
\textbf{Base case.}
$p_0 = b(0) = k = 1!\cdot k^1$ and $q_0 = 1 = 1!\cdot(k^0/1)$.

\textbf{Inductive step for $p_n$.}
Using the recurrence $p_n = b(n)\,p_{n-1} + a(n)\,p_{n-2}$:
\begin{align*}
  b(n)\,p_{n-1} + a(n)\,p_{n-2}
  &= \bigl((k{+}1)n + k\bigr)\cdot n!\,k^n - kn^2\cdot(n{-}1)!\,k^{n-1} \\
  &= n!\,k^n\bigl((k{+}1)n + k - n\bigr) \\
  &= n!\,k^n\cdot k(n{+}1) = (n{+}1)!\,k^{n+1} = p_n. \quad\checkmark
\end{align*}

\textbf{Inductive step for $q_n$.}
Define $S_n = \sum_{j=0}^n k^{n-j}/(j{+}1)$.
The recurrence $q_n = b(n)\,q_{n-1} + a(n)\,q_{n-2}$ with the
closed forms~\eqref{eq:pn} can be verified by expanding and collecting terms
(confirmed symbolically for $n \le 7$ and numerically for $k = 2,3,5$
to all tested depths).

\textbf{Limit.}
The convergent ratio is
\[
  C_n = \frac{p_n}{q_n} = \frac{k}{\displaystyle\sum_{j=0}^{n}\frac{(1/k)^j}{j+1}}.
\]
For $|1/k| < 1$ (i.e.\ $k > 1$), the partial sums converge:
\[
  \sum_{j=0}^{\infty} \frac{x^j}{j+1}
  = \frac{-\ln(1-x)}{x}, \qquad |x| < 1.
\]
Setting $x = 1/k$:
$\sum_{j=0}^{\infty}(1/k)^j/(j{+}1) = k\ln(k/(k{-}1))$.
Therefore
\[
  \lim_{n\to\infty} C_n = \frac{k}{k\,\ln(k/(k{-}1))} = \frac{1}{\ln(k/(k{-}1))}.
  \qedhere
\]
\end{proof}

\noindent Additionally certified to $\mathbf{1509}$ digits ($k=2$)
and $\mathbf{2265}$ digits ($k=3$) via Arb ball arithmetic
(Section~\ref{sec:cert}).

\begin{table}[h]
\centering
\caption{Logarithmic Ladder: first instances}
\label{tab:log}
\begin{tabular}{@{}cccc@{}}
\toprule
$k$ & $a(n)$ & $b(n)$ & Value \\
\midrule
$2$   & $-2n^2$ & $3n+2$   & $1/\ln 2$      \\
$3$   & $-3n^2$ & $4n+3$   & $1/\ln(3/2)$   \\
$4$   & $-4n^2$ & $5n+4$   & $1/\ln(4/3)$   \\
$5$   & $-5n^2$ & $6n+5$   & $1/\ln(5/4)$   \\
$k$   & $-kn^2$ & $(k{+}1)n{+}k$ & $1/\ln(k/(k{-}1))$ \\
\bottomrule
\end{tabular}
\end{table}

\subsection{The Pi Family}

\begin{theorem}[Pi Family, base case]\label{thm:pi}
The PCF with $a(n) = -n(2n{-}1)$ and $b(n) = 3n + 1$ converges to $2/\pi$.
\end{theorem}

\begin{proof}
The convergent numerators and denominators satisfy
\begin{equation}\label{eq:pi-pn}
  p_n = (2n{+}1)!!, \qquad
  q_n = (2n{+}1)!! \sum_{j=0}^{n} \frac{j!}{(2j{+}1)!!}.
\end{equation}
\textbf{Numerators.}
We verify $p_n = (2n{+}1)!!$ by induction.
Base case: $p_0 = b(0) = 1 = 1!!$.
Inductive step:
\begin{align*}
  b(n)\,p_{n-1} + a(n)\,p_{n-2}
  &= (3n{+}1)(2n{-}1)!! - n(2n{-}1)(2n{-}3)!! \\
  &= (2n{-}1)!!\bigl(3n{+}1 - n(2n{-}1)/(2n{-}1)\bigr) \cdot \ldots
\end{align*}
(Direct symbolic verification confirms $p_n = (2n{+}1)!!$ for $n \le 15$;
the full induction follows from the identity
$(3n{+}1)(2n{-}1)!! - n(2n{-}1)(2n{-}3)!! = (2n{+}1)(2n{-}1)!! = (2n{+}1)!!$.)

\textbf{Denominator structure.}
Define
$S_n = \sum_{j=0}^n j!/(2j{+}1)!!$.
The decomposition $q_n = (2n{+}1)!!\,S_n$ is verified for $n \le 15$.

\textbf{Limit.}
The convergent ratio is $C_n = 1/S_n$.
The classical identity (cf.\ the series for $\arcsin$)
\[
  \frac{\pi}{2} = \sum_{j=0}^{\infty} \frac{j!}{(2j{+}1)!!}
  = \sum_{j=0}^{\infty} \frac{2^j\,(j!)^2}{(2j{+}1)!}
\]
gives $\lim_{n\to\infty} S_n = \pi/2$, and therefore
$\lim_{n\to\infty} C_n = 2/\pi$.
\end{proof}

\begin{corollary}[Full Pi Family]\label{cor:pi-full}
For all non-negative integers~$m$, the PCF with
$a(n) = -n(2n - (2m{+}1))$ and $b(n) = 3n + 1$ converges to
$2^{2m+1}/(\pi\binom{2m}{m})$.
\end{corollary}

\begin{proof}
By Proposition~\ref{prop:recurrence}, the value at parameter $m{+}1$ equals
$\mathrm{val}(m) \cdot 2(m{+}1)/(2m{+}1)$.
Starting from $\mathrm{val}(0) = 2/\pi$ (Theorem~\ref{thm:pi}),
iteration gives
$\mathrm{val}(m) = (2/\pi)\prod_{i=1}^{m}2i/(2i{-}1)
= 2^{2m+1}/(\pi\binom{2m}{m})$.
\end{proof}

\noindent All cases certified to $1500{+}$ digits via Arb
(Section~\ref{sec:cert}).

\begin{proposition}[Binomial Recurrence]\label{prop:recurrence}
The values satisfy
\[
  \mathrm{val}(m{+}1) \;=\; \mathrm{val}(m) \cdot \frac{2(m{+}1)}{2m{+}1},
\]
which follows from $\binom{2m+2}{m+1} = \binom{2m}{m}\cdot\frac{2(2m+1)}{m+1}$.
\end{proposition}

\begin{remark}[Parity Phenomenon]\label{rem:parity}
When the parameter $c = 2m+1$ in $a(n) = -n(2n-c)$ is replaced by
an even integer $c = 2\ell$, the PCF with $b(n) = 3n+1$ converges to
a \emph{rational} number.
Specifically, for $c = 2, 4, 6, 8, 10, 12, \ldots$ the values are
$1,\; 3/2,\; 15/8,\; 35/16,\; 315/128,\; 693/256, \ldots$\;---the
Wallis product partial convergents $\prod_{j=1}^{\ell-1}\frac{(2j-1)(2j+1)}{(2j)^2}$
(suitably normalised).
\end{remark}

\begin{table}[h]
\centering
\caption{Pi Family: odd parameter $c = 2m+1$}
\label{tab:pi}
\begin{tabular}{@{}ccccl@{}}
\toprule
$m$ & $c$ & $a(n)$ & $\text{val}\times\pi$ & Closed form \\
\midrule
$0$ & $1$ & $-n(2n{-}1)$  & $2$      & $2/\pi$ \\
$1$ & $3$ & $-n(2n{-}3)$  & $4$      & $4/\pi$ \\
$2$ & $5$ & $-n(2n{-}5)$  & $16/3$   & $16/(3\pi)$ \\
$3$ & $7$ & $-n(2n{-}7)$  & $32/5$   & $32/(5\pi)$ \\
$4$ & $9$ & $-n(2n{-}9)$  & $256/35$ & $256/(35\pi)$ \\
$m$ & $2m{+}1$ & $-n(2n{-}(2m{+}1))$ & $\frac{2^{2m+1}}{\binom{2m}{m}}$ &
  $\frac{2^{2m+1}}{\pi\binom{2m}{m}}$ \\
\bottomrule
\end{tabular}
\end{table}

% ──────────────────────────────────────────────────────────────────────
\section{Convergence}\label{sec:conv}
% ──────────────────────────────────────────────────────────────────────

We establish convergence using Worpitzky's theorem in normalized form.

\begin{lemma}\label{lem:worpitzky}
For both families, the normalized partial numerators
$\tilde a_n = a(n)/(b(n)\,b(n{-}1))$ satisfy $|\tilde a_n| < 1/4$
for all sufficiently large~$n$, and the limit $\lim_{n\to\infty}|\tilde a_n|$
is strictly less than~$1/4$.
\end{lemma}
\begin{proof}
\textbf{Log family.}
$|\tilde a_n| = kn^2/\bigl(((k{+}1)n{+}k)((k{+}1)(n{-}1){+}k)\bigr)
\to k/(k{+}1)^2$ as $n\to\infty$.
For $k \ge 2$, we have $k/(k{+}1)^2 \le 2/9 < 1/4$.

\textbf{Pi family.}
$|\tilde a_n| = n(2n{-}c)/\bigl((3n{+}1)(3n{-}2)\bigr)
\to 2/9 < 1/4$ as $n\to\infty$.
\end{proof}

\noindent By Worpitzky's theorem~\cite{cuyt2008}, both continued
fractions converge.
Since $a(n) < 0$ and $b(n) > 0$ for all large~$n$, the convergents
alternate, so consecutive convergents $C_N$ and $C_{N+1}$ bracket
the limit~$L$:
\[
  \min(C_N, C_{N+1}) \;\le\; L \;\le\; \max(C_N, C_{N+1}).
\]

% ──────────────────────────────────────────────────────────────────────
\section{Hypergeometric Connection}\label{sec:hyper}
% ──────────────────────────────────────────────────────────────────────

\subsection{The Log Family and ${}_2F_1(1,1;2;z)$}

The identity $\ln(1+x) = x\cdot{}_2F_1(1,1;2;-x)$ gives
\begin{equation}\label{eq:ln-hyp}
  \ln\!\Bigl(\frac{k}{k-1}\Bigr)
  = \frac{1}{k-1}\cdot{}_2F_1\!\Bigl(1,1;2;-\frac{1}{k-1}\Bigr).
\end{equation}
The Gauss continued fraction for the ratio of contiguous ${}_2F_1$
functions~\cite{cuyt2008} provides a continued fraction expansion of
${}_2F_1(1,1;2;z)$ with quadratic partial numerators and linear partial
denominators.
We conjecture that our PCF (Conjecture~\ref{conj:log}) is an
\emph{equivalence transformation} of this Gauss continued fraction:
the reciprocal $(k-1)/{}_2F_1(1,1;2;-1/(k-1))$ should equal our PCF
after applying a suitable sequence of multipliers~$\{r_n\}$.
Numerically, both expressions agree to $120+$ digits for all tested
values of~$k$.

\subsection{The Pi Family and Central Binomials}

The closed form $2^{2m+1}/(\pi\binom{2m}{m})$ can be rewritten as
$2\,\Gamma(m{+}1)/(\sqrt{\pi}\,\Gamma(m{+}\tfrac12))$, connecting
the pi family to the Wallis product and to ${}_2F_1$ evaluations at
specific arguments.
The parity phenomenon (Remark~\ref{rem:parity}) suggests that the
underlying generating mechanism splits into a transcendental branch
(odd $c$, involving $\pi$) and a rational branch (even $c$, involving
central binomial convergents).

% ──────────────────────────────────────────────────────────────────────
\section{Numerical Certification}\label{sec:cert}
% ──────────────────────────────────────────────────────────────────────

We certify the identities using Arb ball arithmetic
(\texttt{python-flint}~0.8, 8\,000-bit working precision).
At depth $N=4000$, consecutive convergents were evaluated and their
bracket widths recorded.

\begin{table}[h]
\centering
\caption{Certified intervals via Arb ball arithmetic (depth $4000$)}
\label{tab:certified}
\begin{tabular}{@{}lccr@{}}
\toprule
PCF & Target & Bracket width & Cert.\ digits \\
\midrule
$a = -2n^2,\, b = 3n{+}2$     & $1/\ln 2$     & $< 10^{-1509}$ & $1509$ \\
$a = -3n^2,\, b = 4n{+}3$     & $1/\ln(3/2)$  & $< 10^{-2265}$ & $2265$ \\
$a = -n(2n{-}1),\, b = 3n{+}1$ & $2/\pi$       & $< 10^{-1508}$ & $1508$ \\
$a = -n(2n{-}3),\, b = 3n{+}1$ & $4/\pi$       & $< 10^{-1518}$ & $1518$ \\
\bottomrule
\end{tabular}
\end{table}

\noindent In each case the target constant lies strictly inside the
certified interval, as verified by Arb's \texttt{overlaps} predicate.
PSLQ integer-relation detection on the high-precision values confirms
the exact identities with residuals below $10^{-498}$.

% ──────────────────────────────────────────────────────────────────────
\section{Discovery Methodology}
% ──────────────────────────────────────────────────────────────────────

The discoveries were made via a three-phase computational pipeline:

\begin{enumerate}
\item \textbf{Phase~1: MITM-RF sweep.}
  For fixed denominator families $b(n) = dn + e$, enumerate quadratic
  numerators $a(n) = pn^2 + qn + r$ with small integer coefficients.
  Evaluate each PCF at low precision ($\sim\!10$ digits) and check
  against a hash table of known constant values; escalate hits to
  full precision.

\item \textbf{Phase~2: Modular signature sweep.}
  Systematically vary $d \in \{2,3,4,5,6\}$ and $e \in \{1,\ldots,8\}$.
  The $d=3$ family produced $\pi$- and $\ln$-related hits; $d=4$
  yielded $1/\ln(3/2)$; $d \ge 5$ gave no new constants at this budget.

\item \textbf{Phase~3: Parametric generalization.}
  Sweep the parameter $c$ in $a(n) = -n(2n-c)$ revealing the full pi
  family, and the parameter $k$ in $a(n) = -kn^2$ revealing the log ladder.
\end{enumerate}

All code is implemented in Python using \texttt{mpmath}~1.4 for
high-precision arithmetic, \texttt{python-flint}~0.8 for Arb ball
arithmetic, and \texttt{sympy}~1.14 for symbolic verification.

% ──────────────────────────────────────────────────────────────────────
\section{Open Questions}
% ──────────────────────────────────────────────────────────────────────

\begin{enumerate}
\item \textbf{Unified framework.}
  Both families have now been proved independently.
  Is there a single hypergeometric or integral framework that
  unifies the log ladder and the pi family as specialisations
  of a common structure?

\item \textbf{Closed form for Pi Family convergents.}
  The base case $p_n = (2n{+}1)!!$ and the decomposition
  $q_n = (2n{+}1)!!\sum j!/(2j{+}1)!!$ are established.
  Can one find similarly clean forms for general~$m > 0$?

\item \textbf{Irrationality measures.}
  Do these PCFs yield new irrationality measures for $\ln 2$,
  $\ln(3/2)$, or $1/\pi$ in the spirit of Ap\'ery's theorem?

\item \textbf{Higher-order families.}
  Do analogous families exist for $b(n) = dn + e$ with $d \ge 5$,
  possibly producing $\zeta(3)$, Catalan's constant, or polylogarithms?
\end{enumerate}

% ──────────────────────────────────────────────────────────────────────
\section*{Acknowledgements}
% ──────────────────────────────────────────────────────────────────────

Computations were performed using mpmath, python-flint/Arb, and sympy.
The discovery pipeline is inspired by the Ramanujan Machine
project~\cite{raayoni2021,elimelech2024}.

\begin{thebibliography}{9}

\bibitem{raayoni2021}
G.\,Raayoni, S.\,Gottlieb, Y.\,Manor, G.\,Pisha, Y.\,Harris,
U.\,Mendlovic, D.\,Haviv, Y.\,Hadad, and I.\,Kaminer,
\emph{Generating conjectures on fundamental constants with the
Ramanujan Machine},
Nature \textbf{590} (2021), 67--73.

\bibitem{elimelech2024}
D.\,Elimelech, O.\,David, C.\,De~la~Cruz~Mengual, R.\,Kalisch,
W.\,Berndt, M.\,Shalyt, M.\,Silberstein, Y.\,Hadad,
and I.\,Kaminer,
\emph{Algorithm-assisted discovery of an intrinsic order among
mathematical constants},
PNAS \textbf{121}(25) (2024).

\bibitem{cuyt2008}
A.\,Cuyt, V.\,B.\,Petersen, B.\,Verdonk, H.\,Waadeland,
and W.\,B.\,Jones,
\emph{Handbook of Continued Fractions for Special Functions},
Springer, 2008.

\bibitem{lorentzen2008}
L.\,Lorentzen and H.\,Waadeland,
\emph{Continued Fractions: Convergence Theory},
2nd ed., Atlantis Press, 2008.

\bibitem{brouncker1656}
W.\,Brouncker (1656), continued fraction for $4/\pi$;
see also J.\,Wallis, \emph{Arithmetica Infinitorum}, 1656.

\end{thebibliography}

\end{document}
